% ============================================================================
% 1 INTRODUCTION  —  target: ~0.9 page
% ============================================================================
\section{Introduction}
\label{sec:introduction}

Information extraction (IE) and text classification turn unstructured text
into structured, usable data. A service engineer repairs a machine and records
the customer request, the observed failure, and the corrective action. Across
thousands of reports, these records contain information required for quality
analysis in a form that conventional dashboards cannot directly process. IE
extracts the relevant facts, whereas classification assigns each document to a
category \citep{xu2024generative}.

\begin{figure*}[!t]
\centering
\resizebox{\textwidth}{!}{%
\begin{tikzpicture}[node distance=4mm]
% ---- Row 1: classical supervised ML -------------------------------------
\node[tumdata] (c1) {Large ticket set};
\node[tumbox, right=7mm of c1] (c2) {Labeling by\\domain expert\\\scriptsize slow and expensive};
\node[tumdata, right=7mm of c2] (c3) {Train / test split};
\node[tumbox, right=7mm of c3] (c4) {Train and tune\\\scriptsize by an ML expert};
\node[tumdata, right=7mm of c4] (c5) {Inference};
\foreach \a/\b in {c1/c2, c2/c3, c3/c4, c4/c5} \draw[tumarrow] (\a) -- (\b);
\node[tumlabel, left=5mm of c1, align=right, anchor=east] {\textbf{(a) Classical ML}};

% ---- Row 2: LLM only ----------------------------------------------------
\node[tumdata, below=9mm of c1] (l1) {Large ticket set};
\node[tumbox, right=7mm of l1] (l2) {LLM};
\node[tumdata, right=7mm of l2] (l3) {Inference};
\foreach \a/\b in {l1/l2, l2/l3} \draw[tumarrow] (\a) -- (\b);
\node[tumlabel, left=5mm of l1, align=right, anchor=east] {\textbf{(b) LLM only}};
\node[tumlabel, right=7mm of l3, text=TUMOrange, align=left, anchor=west]
  {\;$\leftarrow$ no step here measures how well it works};

% ---- Row 3: LLM with evaluation (what we built) -------------------------
\node[tumdata, below=9mm of l1] (e1) {Large ticket set};
\node[tumbox2, right=7mm of e1] (e2) {\textbf{Step 1}\\LLM pre-annotation\\\scriptsize spans, labels, offsets};
\node[tumbox2, right=7mm of e2] (e3) {\textbf{Step 2}\\Expert validation\\\scriptsize confirm $\cdot$ correct $\cdot$ delete};
\node[tumdata, right=7mm of e3] (e4) {Gold test set};
\node[tumbox, right=7mm of e4] (e5) {Evaluation\\\scriptsize P / R / F$_1$};
\node[tumdata, right=7mm of e5] (e6) {Inference};
\foreach \a/\b in {e1/e2, e2/e3, e3/e4, e4/e5, e5/e6} \draw[tumarrow] (\a) -- (\b);
\node[tumlabel, left=5mm of e1, align=right, anchor=east] {\textbf{(c) LLM with evaluation}};

\begin{scope}[on background layer]
\node[draw=TUMOrange!50, dashed, rounded corners=3pt, inner sep=6pt, fit=(e2)(e3)] (tool) {};
\node[draw=TUMBlue!40, dashed, rounded corners=3pt, inner sep=6pt, fit=(e5)] (fw) {};
\end{scope}
\node[tumlabel, below=1mm of tool.south, anchor=north, text=TUMOrange] {annotation tool};
\node[tumlabel, below=1mm of fw.south, anchor=north, text=TUMBlue] {evaluation framework};
\end{tikzpicture}%
}
\caption{Three annotation workflows. (a)~The supervised pipeline requires
expert labeling before model development. (b)~Direct LLM inference omits an
independently labeled evaluation set. (c)~The implemented workflow combines
LLM pre-annotation with expert validation and gold-set evaluation. Dashed
boxes denote the two artifacts developed in this project.}
\label{fig:workflows}
\end{figure*}

Until recently, this process required substantial initial investment,
workflow~(a) in \Cref{fig:workflows}: domain experts labeled thousands of
examples, machine-learning specialists trained and tuned a model, and label-set
changes could require additional annotation and training. Instruction-tuned
large language models (LLMs) reduce this initial effort because a task can be
specified through label descriptions and examples in a prompt, workflow~(b).

The principal limitation is the loss of independent evaluation data. The
labels created for supervised learning also served as a test set against which
predictions could be measured. Organizations may therefore lack annotations
for their target domain and cannot reliably evaluate deployed models. Vendor
benchmarks may differ substantially from service reports, while informal spot
checks cannot establish deployment quality. LLM annotation quality varies
across tasks and datasets and must therefore be measured against human labels
rather than assumed
\citep{pangakis2023validation, reiss2023testing}.

Workflow~(c) is implemented in this study. An LLM pre-annotates, a domain expert
validates, and the result is a small gold set to measure against. We built both
halves: an evaluation framework for classification and span extraction, and a
configurable annotation product that turns model output into reviewed data.
The initial Few-NERD and support-ticket experiments test method and output
choices on short inputs. The later SciREX study addresses the harder operational
case: coherent scientific text ranging from one sentence to more than 200,
processed completely with durable checkpoints and evaluated against released
span annotations \citep{jain2020scirex}. This progression matters because
sentence-local predictions that work on isolated inputs must still be shown to
cover every sentence in a full document reliably and economically.

\paragraph{Contributions.}
Our contributions are threefold. First, we provide a task-separated,
configuration-driven evaluation framework that stores raw responses and scores
them independently of inference. Second, we construct a reproducible
1{,}000-window SciREX benchmark from 438 source papers while preserving source
order, document identity, split membership, and exact character offsets; the
completed study stages include a development prompt comparison, a fresh
20-paper held-out confirmation, and a 100-paper full-window operational run.
Third, we deliver a Dockerized annotation application with schema-derived but
editable prompts, arbitrary labels, text/PDF/CSV input, optional reviewer-
approved LLM or local-embedding ticket routing, confidence-ranked review, and
auditable JSON export.
Together, these artifacts define a reproducible process from unlabeled
organizational text to an evaluation set with inspectable provenance.
